\section{Results}
\label{sec:results}

We organize the results around the offline architectural question:
whether pair descriptors should affect only attention weights or also
the aggregated value messages, and how this choice interacts with bias
sharing. Additional relation features are a secondary extension of that
comparison. The official full-test metrics remain pending; current
offline numbers are explicitly marked drafting placeholders from the
earlier 500,000-jet evaluation.

Multi-seed summaries average metrics from separately evaluated checkpoints,
not ensemble predictions. Quoted seed spreads are sample standard
deviations across three matched training seeds. Macro AUC is the equally
weighted ten-class one-versus-rest mean, including QCD, on the 0--1 scale.
We summarize the earlier reduced-view evidence after the offline results;
its complete methods and measurements are retained in
Appendix~\ref{sec:hlt-support}.

\subsection{Offline results on the full test set}
\label{sec:offline-full-test-results}

\begin{quote}
\textbf{Draft placeholders---20M results pending.} Every numerical offline
result below, including derived comparisons, seed details and rejection
tables, is temporarily carried over from the earlier 500,000-jet test
(50,000 QCD jets). A dagger marks borrowed table entries; TBD denotes an
unavailable ablation. These values and the accompanying interpretation must
be updated from the full-test artifacts before publication. They are not
20M-test measurements.
\end{quote}

The primary offline comparison uses the official 20-million-jet test set,
with two million jets per class. Training remains fixed at one million jets
per model and three seeds using the same training sample. Thus, the larger
evaluation tests the frozen models more precisely rather than increasing
their training data. Table~\ref{tab:offline-full-test} organizes the seven
configurations: four have trained checkpoints undergoing full-test
evaluation, while the three additional ablations still require training
and evaluation.

% Draft 20M table: numerical placeholders come from the original 500k test.
% Replace every provisional cell from authenticated 20M artifacts.
% AUC is first averaged over classes per seed, then over seeds.
\begin{table}[htbp]
\centering
\footnotesize
\setlength{\tabcolsep}{3pt}
\begin{tabularx}{\linewidth}{@{}Xrrr@{}}
\toprule
Configuration & Accuracy (\%) & $\Delta$ (pp) & Macro AUC \\
\midrule
Standard: shared & \provisionalresult{$82.8519 \pm 0.0982$} & --- & \provisionalresult{$0.982317 \pm 0.000190$} \\ % data:offline_placeholder_500k:OFF_RPT_BASE
Standard: layerwise & \pendingresult & \pendingresult & \pendingresult \\
Standard: shared + EV & \pendingresult & \pendingresult & \pendingresult \\
Standard: layerwise + EV & \provisionalresult{$83.4722 \pm 0.0510$} & \provisionalresult{+0.6203} & \provisionalresult{$0.983325 \pm 0.000039$} \\ % data:offline_placeholder_500k:OFF_RPT_BASE_EDGEVALUE
Selected: layerwise & \provisionalresult{$83.5152 \pm 0.0885$} & \provisionalresult{+0.6633} & \provisionalresult{$0.983464 \pm 0.000118$} \\ % data:offline_placeholder_500k:OFF_RPT_SELECTED_LAYERWISE
Selected: shared + EV & \pendingresult & \pendingresult & \pendingresult \\
Selected: layerwise + EV & \provisionalresult{$83.8670 \pm 0.0385$} & \provisionalresult{+1.0151} & \provisionalresult{$0.984010 \pm 0.000130$} \\ % data:offline_placeholder_500k:OFF_RPT_SELECTED_EDGEVALUE
\bottomrule
\end{tabularx}
\caption{Offline comparison intended for the official 20M test.
\textbf{Draft: daggered values are 500k-test placeholders, not 20M results.}
TBD denotes the three additional, untrained ablations. All models use one
million training jets and three seeds. Accuracy and macro AUC are means
with sample standard deviations; AUC uses the 0--1 scale and differences
are relative to the standard shared-bias baseline. Selected means standard
features plus PT, TRACK, and REGION; EV denotes EdgeValue.}
\label{tab:offline-full-test}
\end{table}


\paragraph{Separating value messages from bias sharing.}
The most direct available value-path comparison holds the selected
relations and layerwise bias fixed. Adding EdgeValue increases mean
accuracy by 0.3518 percentage points in the provisional measurements,
with a positive difference at each of the three matched seeds. This
supports the value pathway within the tested architecture, but does not
establish that layerwise bias is necessary for it.

With standard pair features, layerwise bias plus EdgeValue improves
accuracy over the shared-bias baseline by 0.6203 percentage points in the
placeholder values. Because both bias sharing and the value pathway
change, that difference cannot be assigned to EdgeValue alone. The
standard-layerwise-only and standard-shared-bias EdgeValue ablations
are pending. They are needed to complete the feature-controlled
$2\times2$ architectural comparison and test how the two choices interact.

\paragraph{Combined design and additional relation features.}
The provisional selected-relation, layerwise-bias EdgeValue model reaches
a mean accuracy of 83.8670\% (seed SD 0.0385 percentage points), compared
with 82.8519\% (SD 0.0982) for the standard baseline. Its advantage of
1.0151 percentage points corresponds to a 5.92\% relative reduction in mean
classification error. This is the gain of the combined design, not a
standalone estimate of the EdgeValue contribution. Macro AUC increases
from $0.982317\pm0.000190$ to $0.984010\pm0.000130$; the standard-feature
EdgeValue and selected layerwise-only models reach
$0.983325\pm0.000039$ and $0.983464\pm0.000118$, respectively.

The selected EdgeValue model exceeds the other three available
configurations at every matched seed in these provisional measurements.
Its accuracy gains over baseline range from 0.9276 to 1.1710 percentage
points. The selected layerwise-only model improves over baseline by
0.6633 percentage points, but that comparison changes both pair features
and bias sharing. The selected shared-bias EdgeValue ablation will extend
the comparison of bias sharing to the augmented feature setting.
Seed-level placeholders and paired-interval details appear in
Appendix~\ref{sec:offline-details}. Full-test results must establish
whether the provisional gains and ordering persist; the present evidence
does not establish a universally optimal feature set or gains at larger
training-set sizes.

\paragraph{QCD background rejection.}
Tables~\ref{tab:offline-rejection-30} and~\ref{tab:offline-rejection-50}
reserve the classwise comparison at 30\% and 50\% signal efficiency. They
currently contain placeholders for the four trained configurations; the
three untrained ablations have no rejection measurements. With the
placeholder values, selected EdgeValue exceeds the baseline at every
tabulated operating point where both means are finite. At 50\% efficiency,
the increases range from 2.63\% for $W\to q\bar q'$ to 87.15\% for
$H\to b\bar b$. At 30\%, the finite increases range from 6.03\% for
$Z\to q\bar q$ to 66.89\% for $H\to4q$.

% Draft 20M presentation; all entries are placeholders from the 500k test.
% SAT means at least one zero-background seed in that source sample.
\begin{table}[htbp]
\centering
\footnotesize
\setlength{\tabcolsep}{2.5pt}
\begin{tabularx}{\linewidth}{@{}Xrrrrr@{}}
\toprule
Signal & \provisionalresult{B} & \provisionalresult{BE} & \provisionalresult{SL} & \provisionalresult{SE} & \provisionalresult{SE/B change} \\
\midrule
$H\to4q$ & 3582.82 & 5218.86 & 6613.76 & 5979.44 & +66.89\% \\ % data:rejection_30:H4q
$H\to b\bar b$ & 22222.22 & 41666.67 & 41666.67 & SAT & --- \\ % data:rejection_30:Hbb
$H\to c\bar c$ & 4444.44 & 7010.58 & 4557.48 & 6296.30 & +41.67\% \\ % data:rejection_30:Hcc
$H\to gg$ & 315.16 & 339.29 & 343.26 & 359.44 & +14.05\% \\ % data:rejection_30:Hgg
$H\to q\bar q\ell\nu$ & SAT & SAT & SAT & SAT & --- \\ % data:rejection_30:Hqql
$t\to b\ell\nu$ & SAT & SAT & SAT & SAT & --- \\ % data:rejection_30:Tbl
$t\to bq\bar q'$ & SAT & SAT & SAT & SAT & --- \\ % data:rejection_30:Tbqq
$W\to q\bar q'$ & 1575.31 & 1598.52 & 1857.86 & 2045.86 & +29.87\% \\ % data:rejection_30:Wqq
$Z\to q\bar q$ & 1076.71 & 1220.00 & 1345.17 & 1141.65 & +6.03\% \\ % data:rejection_30:Zqq
\bottomrule
\end{tabularx}
\caption{Offline mean QCD rejection at 30\% signal efficiency.
\textbf{Draft: all entries under daggered headers are 500k-test placeholders,
not 20M results; their background support is 50,000 QCD jets.}
B is the baseline; BE uses standard
features, layerwise bias and EdgeValue; SL uses selected features and
layerwise bias; SE additionally uses EdgeValue. Means average the three
per-seed rejections, not pooled passing counts. SAT denotes no finite
three-seed mean because at least one seed passes zero QCD events.
Percentage changes are ratios of means calculated before rounding.
The three untrained ablations have no entries yet.}
\label{tab:offline-rejection-30}
\end{table}

% Draft 20M presentation; all entries are placeholders from the 500k test.
% SAT means at least one zero-background seed in that source sample.
\begin{table}[htbp]
\centering
\footnotesize
\setlength{\tabcolsep}{2.5pt}
\begin{tabularx}{\linewidth}{@{}Xrrrrr@{}}
\toprule
Signal & \provisionalresult{B} & \provisionalresult{BE} & \provisionalresult{SL} & \provisionalresult{SE} & \provisionalresult{SE/B change} \\
\midrule
$H\to4q$ & 750.41 & 791.42 & 859.64 & 1001.51 & +33.46\% \\ % data:rejection_50:H4q
$H\to b\bar b$ & 4848.48 & 4788.96 & 8928.57 & 9074.07 & +87.15\% \\ % data:rejection_50:Hbb
$H\to c\bar c$ & 1411.91 & 1360.31 & 1548.40 & 1482.20 & +4.98\% \\ % data:rejection_50:Hcc
$H\to gg$ & 96.28 & 98.18 & 98.19 & 101.20 & +5.12\% \\ % data:rejection_50:Hgg
$H\to q\bar q\ell\nu$ & SAT & SAT & SAT & SAT & --- \\ % data:rejection_50:Hqql
$t\to b\ell\nu$ & SAT & SAT & SAT & SAT & --- \\ % data:rejection_50:Tbl
$t\to bq\bar q'$ & 5092.59 & 7566.14 & 6845.24 & 7229.44 & +41.96\% \\ % data:rejection_50:Tbqq
$W\to q\bar q'$ & 320.55 & 316.53 & 322.22 & 328.98 & +2.63\% \\ % data:rejection_50:Wqq
$Z\to q\bar q$ & 236.62 & 267.42 & 260.29 & 265.12 & +12.04\% \\ % data:rejection_50:Zqq
\bottomrule
\end{tabularx}
\caption{Offline mean QCD rejection at 50\% signal efficiency.
\textbf{Draft: all entries under daggered headers are 500k-test placeholders,
not 20M results; their background support is 50,000 QCD jets.}
B is the baseline; BE uses standard
features, layerwise bias and EdgeValue; SL uses selected features and
layerwise bias; SE additionally uses EdgeValue. Means average the three
per-seed rejections, not pooled passing counts. SAT denotes no finite
three-seed mean because at least one seed passes zero QCD events.
Percentage changes are ratios of means calculated before rounding.
The three untrained ablations have no entries yet.}
\label{tab:offline-rejection-50}
\end{table}


The provisional rejection ordering is not uniform across architectures and
classes. Selected layerwise bias without EdgeValue has higher rejection
than selected EdgeValue for $H\to4q$ and $Z\to q\bar q$ at 30\%
efficiency, and for $H\to c\bar c$ at 50\%. Higher aggregate accuracy
therefore does not imply the best result at every class-specific operating
point.

The small background support of the placeholder measurements is especially
important in the tails. For $H\to b\bar b$ at 50\% efficiency, their
baseline passing counts are 10, 10, and 11 QCD jets, versus 9, 10, and 3
for selected EdgeValue. The large rejection increase is consequently
sensitive to one low-count seed. SAT marks an operating point with zero
passing QCD jets for at least one seed, not infinite population rejection.
The full test supplies two million QCD jets rather than the placeholders'
50,000; all rejection values, saturation flags, counts, achieved
efficiencies and percentage changes must be recomputed for that sample.
Full-test conditional intervals, as defined in
Section~\ref{sec:evaluation-protocol}, are pending and cannot be inferred
from these placeholders.

\subsection{Supporting evidence from the synthetic reduced-input study}
\label{sec:hlt-support-summary}

The earlier HLT-like study provides exploratory context, rather than a
second headline benchmark. Its 21-configuration validation screen selected
PT+TRACK+REGION, fixing the additional pair descriptors before the offline
campaign. This documents the origin of the selected feature set without
claiming that a broad offline relation-family search has been completed.

The three-seed architecture comparisons in that reduced view also motivate
testing the value pathway separately from bias sharing. Adding EdgeValue
to an already-layerwise model improves mean test accuracy by 0.3089
percentage points with standard features and 0.2633 with selected
features. Layerwise bias alone gives smaller and feature-dependent gains
of 0.0175 and 0.1987 percentage points, respectively; its standard-feature
macro AUC is slightly below the shared-bias baseline. The supporting study
therefore does not establish a universal benefit from layerwise bias or
that it is required for EdgeValue.

These comparisons use a fixed, hand-designed input degradation, not
validated detector-level trigger reconstruction, and share underlying jet
identities with the original offline development samples. They do not
constitute independent-dataset replication. The complete appendix retains
all screening and final-test configurations, capacity and unary controls,
and rejection results, including unfavorable comparisons.
